\renewcommand{\MayorVer}{2}
\renewcommand{\MenorVer}{1}
\fancypagestyle{Fext}{%
    \fancyhead[C]{%
	    \begin{tblr}{%	
            colspec = {|X[c,m]|X[c,m]|X[c,m]|X[c,m]|X[c,m]|X[c,m]|},
            }
            \hline
            \SetCell[c=2]{c} \includegraphics[height=1cm]{./RDF_Logo.eps} & & \SetCell[c=3]{c} \large{Red de fríos S.A. de C.V.} & & & {\textbf{Código:}\\ \Codigo} \\
            \hline
            \SetCell[c=4]{c} {{Política de trabajo en caso de personal con heridas expuestas\\ o enfermedades infecciosas}} & & & & {\textbf{Fecha de publicación:}\\ \FechaPub} & {\textbf{Edición:}\\ \Edit} \\
            \hline
            \end{tblr}
	    }%
        \fancyfoot[CE,CO,LE,LO,RE,RO]{} %% clear out all footers
        \fancyfoot[C]{Página \thepage\ de \pageref{LastPage}}
        \fancyfoot[L]{\textbf{CONFIDENCIAL}}
        \fancyfoot[R]{\Codigo~rev.\Edit}
}%

\renewcommand{\Codigo}{BPD-8-POL/IT}
\renewcommand{\FechaPub}{2023--01}


\pagestyle{Fext}

%\section{Política de trabajo en caso de personal con heridas expuestas o enfermedades infecciosas}

\subsection{Objetivo}

\begin{itemize}
	\item \textbf{Evitar} que el personal que padezca alguna enfermedad infecta contagiosa, afecciones en la piel a alguna herida entre en contacto con el producto, para evitar algún tipo de contaminación.
\end{itemize}

\subsection{Alcance}

A todas las personas que presente algún tipo enfermedad infecta contagiosa, afecciones en la piel a alguna herida.

\subsection{Procedimiento}

\subsubsection{Personal con heridas}

\begin{itemize}
	\item Las cortadas y heridas deben de cubrirse apropiadamente con material impermeable, evitando entrar en el \emph{área de almacenamiento} cuando estas estén en contacto directo con el producto y puedan propiciar la contaminación del mismo.
	\item El personal en esta situación debe ser ubicado en áreas donde realice actividades que no tengan contacto con el producto, materia prima o material de empaque (Fuera de áreas operativas) para evitar la contaminación del producto.
	\item Se debe de informar al personal del área operativa que algún trabjador con heridas entrará a trabajar, para que ellos puedan prestar atención a que el trabajador en cuestión no entre en contacto con producto alimenticio.
\end{itemize}

\subsubsection{Personal con enfermedades infecto-contagiosas o afecciones en la piel}

\begin{itemize}
	\item El personal debe de avisar en caso de presentar alguna anormalidad en la piel o la presencia de una enfermedad contagiosa;
	\item se debe evitar que las personas con enfermedades contagiosas laboren en contacto directo con los productos, materia prima, o material de empaque (si su situación lo permite).
	\begin{itemize}
		\item si la enfermedad es grave, el personal no debe de presentarse a su área de trabajo, para evitar cualquier riesgo de contaminación del producto o del personal que labora en la empresa.
	\end{itemize}
	\item El personal debe ser enviado al médico para su revisión.
	\item El personal debe presentarse con el certificado que avale que es apto para trabajar.
	\begin{itemize}
		\item Los productos que se manejan en el almacén no están expuestos ya que cuentan con empaque primario y secundario.
		\item El uniforme del personal cubre en su mayoría el cuerpo del personal, exceptuando los ojos.
	\end{itemize}
\end{itemize}

\subsection{Frecuencia}

Cuando se presente alguna persona enferma.

\subsection{Responsables}
\index{Responsables!Personal con heridas expuestas}
\begin{itemize}
	\item El personal con heridas que \textbf{tenga} que laborar en el área de almacenamiento, deberá ser supervisado en todo momento por el resto del personal de \textbf{operaciones.}
\end{itemize}

\subsection{Acciones preventivas}
\index{Acciones preventivas!Personal con heridas expuestas}
\begin{itemize}
	\item En caso de que un operador con heridas tenga que laborar, se les informará al área correspondiente para que presten atención de que el operador en cuestión nunca entre en contacto con producto expuesto.
\end{itemize}

\subsection{Acciones correctivas}
\index{Acciones correctivas!Personal con heridas expuestas}
\begin{itemize}
	\item En caso de que el personal con heridas entre en contacto directo con el producto expuesto, este tendrá que segregarse e informarse al cliente sobre este suceso;
	\item se llenara el formulario de acciones correctivas y se hará un análisis causa raíz para determinar cambios en el procedimiento correspondiente.
\end{itemize}

\begin{changelog}[simple, sectioncmd=\subsection*,label=changelog-1.8]
	\begin{version}[v=2.1, date=2023--01, author=Pablo E. Alanis]
		%\fixed
			\item Cambio de formato;
			\item se adoptaron versiones mayores y menores para revisiones;
			\item cambio en el código del documento.
	\end{version}

	\begin{version}[v=1.8, date=2022-05, author=Alonso M.]
		\item \textit{No hubo cambios}
	\end{version}
	\shortversion{v=1.7, date=2021-05, changes=No hubo cambios}
\end{changelog}

\pagestyle{formato}